% Chapter 4 session close, 2026-06-18
% Append-only commentary file

\section*{Chapter 4 - session close, 2026-06-18}

Chapter 4 drafted E2E, \S1 through \S6. Six sections, seven paragraphs, plumbing is down but then numbers not yet—this is on purpose as this is a pilot draft.

\subsection*{ThingsToDo}

Three gaps all in \$4, all waiting on the real PCA run. The spectrum and $k$ in \P4. The loadings and dimension names in \P5. The prior-versus-result comparison in \P6. Every framing sentence is written and only the values are bracketed. When the run lands I swap the brackets.

The figures wait on the same data: the spectrum-with-gaps in \S4\,\P4, the loading table is \S4\,\P5. One more table waits on variable request: the candidate indicators in \S2\,\P3. I did not commit to seventeen, or to any number---the count is bracketed because the request is not back (writing now would be fabricating).

Conlon---verify Friday. Source is real but not in tracker, I pulled it from memory. The sentence is written as if the citation is coming. 

The introduction now owes two debts the body spends. It has to establish deprivation, age, and isolation as the field's recurring themes (\S4\,\P3, Klinenberg and several more), and it has to establish age as a leading modifier of heat--mortality (\S6\,\P2, and vary the source---Ch3 already leans on Gasp 2022 for this). The body states both flat and uncited because the introduction is meant to have earned them. That is the whole reason the introduction is written last.

\subsection*{Pointers --- the reasoning, kept so I don't ask ``why did I write this'' later}

Keep-count and name-count are two different budgets. The Akaike criterion decides how many components to carry into the model; the eigenvalue gaps decide how many of those may be named. I can keep four and name two. \S5\,\P4 holds them apart but keeps it quiet---flagged whether the close needs an explicit ``the model may carry more than we name,`` or whether saying it out loud there would over-explain at the landing.

The keep-rule is rendered as AIC throughout, to match Ch2. Settled, unless the real run gives me a reason to open. 

\S6\,\P3 hands to Ch5 by the name of its job, not its number---``the next chapter,`` no numbered reference. This breaks from Ch3, which numbered the forward-reference. Deliberate, per the brief: name the door by what is behind it. Revisit at the full-manuscript pass if the inconsistency grates.

\subsection*{$\rightarrow$ Anki --- word-atoms and moves to harvest Friday}
 
The semicolon is one idea seen twice; the period is two facts. Use the
semicolon when the second clause is the first understood at a deeper
level.
 
The contrapositive defends a design choice: don't argue your choice is
good, show the alternative breaks at the worst possible spot.
 
A reused example is worse than a fresh one --- the reader feels the
repeat, and the point was self-evident the second time anyway.
 
After a clean two-case split, the summary sentence that says ``therefore
it fails'' is redundant; the split already contains the conclusion, and
the reader has drawn it.
 
A list before a short flat turn only works if the list is short enough
that the reader still has the energy to feel the turn. Six items and the
turn is relief; three and it snaps.